% ============================================================
% Appendice C — Template pronti all'uso
% ============================================================

\chapter{Template Pronti all'Uso}
\label{app:template}

Template copiabili per i file più usati di AI-DLC. Sostituisci i valori tra \texttt{[parentesi]} con quelli reali del tuo progetto. I template completi in formato Markdown sono in \texttt{.ai-dlc/modules/templates/}.

\section*{Catalogo degli artefatti AI-DLC}

Il framework produce un insieme definito di documenti. Ognuno ha uno scopo, un momento in cui nasce e un posto dove vive (versionato nel repository). Questa è la mappa; i template seguono.

\begin{table}[htbp]
  \centering\small
  \begin{tabular}{p{0.21\textwidth}p{0.45\textwidth}p{0.24\textwidth}}
    \toprule
    \textbf{Artefatto} & \textbf{Cos'è / quando} & \textbf{Dove vive} \\
    \midrule
    \texttt{\_CONTEXT.md} & Stato corrente + vincoli invalicabili (fonte di verità); ogni sessione & radice progetto \\
    \texttt{PROGRESS.md} & Diario cronologico di sessioni, decisioni, scoperte; in append & radice progetto \\
    Requisiti FR/NFR & Requisiti funzionali e non funzionali; Fase 1 & \texttt{docs/01\_ANALYSIS/} \\
    Glossario dominio & Termini condivisi del dominio; Fase 0--1 & \texttt{docs/01\_ANALYSIS/} \\
    User Story & Requisito lato utente + criteri di accettazione; Fase 1 & backlog \\
    Modello dati & Entità e relazioni, tecnologia-agnostico; Fase 1 & \texttt{docs/01\_ANALYSIS/} \\
    Threat model & Asset, attori, superfici, mitigazioni; Fase 1--2 & \texttt{docs/02\_DATA\_GOVERNANCE/} \\
    ADR & Decisione architetturale: contesto, opzioni, scelta, conseguenze; su decisione HIGH/CRITICAL & \texttt{docs/} (vicino al codice) \\
    EPIC & Contenitore di valore con user story, AC e task; pianificazione feature ampia & backlog / \texttt{docs/} \\
    TASK & Unità di implementazione (sizing, AC, test); Fase 3, prima del codice & task docs \\
    Analisi codebase & MAP/RISKS/HOTSPOTS/RUN di codice esistente; prima di modifiche legacy & \texttt{docs/} \\
    Company process & Gate e standard aziendali; setup enterprise & \texttt{.ai-dlc/company/} \\
    Project instructions & Regole specifiche di progetto; setup & \texttt{.ai-dlc/project/} \\
    \bottomrule
  \end{tabular}
  \caption{Gli artefatti prodotti dal framework AI-DLC}
  \label{tab:appc-artefatti}
\end{table}

\section*{Template 1 --- \texttt{\_CONTEXT.md} Minimo (spike/POC)}

\begin{lstlisting}[language=,caption={Context minimale per spike}]
# PROJECT CONTEXT — MINIMAL
> Ultimo aggiornamento: [YYYY-MM-DD]

| Param                   | Value                              |
|-------------------------|------------------------------------|
| Project                 | [Nome]                             |
| Phase                   | [fase]                             |
| Mode                    | RAPID                              |
| Active Task             | [descrizione task]                 |
| Active Task Token Est.  | 0/0/0                              |
| Active Task Model Level | 3                                  |

## Stack
| Runtime   | [tecnologia]  |
| Framework | [framework]   |

## Vincoli Essenziali
| SEC-01 | Input Validation | [approccio] |
| SEC-02 | Auth             | [meccanismo] |

## Note per l'Agente
- Mode RAPID: branch dedicato, non toccare produzione
- Obiettivo time-boxed: [cosa stai esplorando, entro quando]
\end{lstlisting}

\section*{Template 2 --- Voce \texttt{PROGRESS.md}}

\begin{lstlisting}[language=,caption={Voce di sessione per PROGRESS.md}]
## [YYYY-MM-DD] — [ID Task] [Titolo Task]

**Fatto:**
- [cosa è stato implementato]

**Scoperto:**
- [comportamento inaspettato, bug latente trovato]

**Decisione:**
[DECISIONE: descrizione]
Motivazione: [perché questa opzione]

**Prossima sessione:** [cosa fare]

---
\end{lstlisting}

\section*{Template 3 --- Task AI-DLC}

\begin{lstlisting}[language=,caption={Template di task con AI Sizing}]
# [TASK-ID] — [Titolo del Task]

## AI Sizing
| Token Estimate     | [input] / [output] / [totale] |
| Model Level        | [1-7]                         |
| Risk Floor Applied | [LOW/MEDIUM/HIGH → min N]     |
| Rationale          | [motivazione]                 |

## Goal
[Cosa deve produrre questo task — 1-3 frasi]

## Acceptance Criteria
- [ ] [criterio verificabile 1]
- [ ] [criterio verificabile 2]

## Test Plan (TDD)
Test da scrivere PRIMA del codice (codificano gli AC):
- [ ] [test caso nominale]
- [ ] [test caso d'errore]
- [ ] [test edge case / sicurezza]
Definizione di "fatto": test del task verdi
+ regressione completa verde.

## Vincoli rilevanti
- [SEC-XX o PERF-XX applicabili]
\end{lstlisting}

\section*{Template 4 --- Override HALT Trigger di Progetto}

\begin{lstlisting}[language=yaml,caption={Override di progetto per halt-triggers.yaml}]
version: 1

# Per disattivare un trigger predefinito (temporaneamente):
# disable:
#   - schema

triggers:
  - id: billing
    patterns:
      - "**/billing/**"
      - "**/payments/**"
    reason: Logica di pagamento — impatto finanziario
    risk: HIGH+
\end{lstlisting}

\section*{Template 5 --- Skill di Progetto Personalizzata}

\begin{lstlisting}[language=,caption={Template di skill personalizzata}]
# SKILL — [Nome Skill]

> Carica quando: [tipo di lavoro]

## Identità
Sei [ruolo specializzato]. Priorità: [principio].

## Regole di Dominio
- [regola di business 1]
- [invariante critica]

## Anti-pattern
- NON fare [cosa] — [perché]

## Checklist di Verifica
- [ ] [controllo 1]
\end{lstlisting}

\section*{Template 6 --- EPIC}

\begin{lstlisting}[language=,caption={Template di EPIC}]
# EPIC: [E-NNN] [Titolo]
Status: Proposed | Active | Done | Deferred
Owner: [team/persona]   Target: [sprint/release]

## Summary
[Valore di business in 2-3 frasi]

## Scope
In scope: [...]
Out of scope: [...]

## User Stories
- [US-001] Come [ruolo] voglio [capacità]
  così da [risultato].

## Acceptance Criteria
- [ ] [AC a livello epic]

## Task Breakdown
| ID | Task | Tipo | Token Est. | Model Lvl | Risk |

## Definition of Done
- [ ] Tutti i task collegati completati
- [ ] Test verdi (task + regressione)
- [ ] SEC/PERF verificati; doc aggiornata
\end{lstlisting}

\section*{Template 7 --- ADR (Architecture Decision Record)}

\begin{lstlisting}[language=,caption={Template di ADR / Decision Record}]
# ADR-NNN: [Titolo decisione]
Status: Proposed | Accepted | Deprecated | Superseded
Date: YYYY-MM-DD   Owner: [team/persona]
Risk: LOW|MEDIUM|HIGH|HIGH+|CRITICAL
Model Level: [1-7]

## Contesto
[Problema, vincoli, riferimenti FR/NFR/SEC/PERF]

## Opzioni
| Opzione | Descrizione | Pro | Contro | Rischio |
| A       | ...         | ... | ...    | ...     |
| B       | ...         | ... | ...    | ...     |

## Decisione
[Opzione scelta e motivazione concisa]

## Conseguenze
Positive: [...]   Negative: [...]
Mitigazioni: [...]

## Follow-up
- [ ] [task/riferimento]
\end{lstlisting}

\section*{Template 8 --- Requisiti (Fase Analysis)}

\begin{lstlisting}[language=,caption={Requisiti funzionali, user story e NFR}]
## Requisiti Funzionali
| ID   | Descrizione | Priorità         | Fonte |
| FR01 | ...         | Must/Should/Could | ...   |

## User Story
[US-XXX] Titolo
Come [ruolo] voglio [azione] così da [beneficio]
Acceptance Criteria:
- [ ] GIVEN [contesto] WHEN [azione] THEN [risultato]
Security: [...]  Performance: [...]
Priorità: Must/Should/Could

## Requisiti Non Funzionali (NFR)
- Sicurezza: SEC-01..SEC-09 applicabili
- Performance: PERF-01..PERF-05 (P95, throughput...)
\end{lstlisting}

\section*{Template 9 --- Analisi di codebase esistente (legacy)}

Prodotti dal modulo \texttt{09\_CODEBASE\_ANALYSIS.md} prima di modificare codice esistente: \texttt{MAP.md} (moduli, dipendenze, entry point), \texttt{RISKS.md} (rischi sicurezza/performance e debito), \texttt{HOTSPOTS.md} (file/classi critici con motivazione), \texttt{RUN.md} (comandi build/test/run).

\begin{lstlisting}[language=,caption={Architecture Snapshot e Change Map}]
## Architecture Snapshot
- Pattern: [layered/clean/hex/monolite modulare]
- Entry point: [...]
- Domini/moduli core: [...]
- Data store: [...]   Integrazioni: [...]

## Change Map (per una modifica specifica)
Target: [feature/bug]
- Percorso primario: [file/classe/funzione]
- Impatti secondari: [...]
- Test da toccare: [...]
- Piano di rollback: [...]
\end{lstlisting}
